\small
\begin{threeparttable}
\begin{tabular}{lccc}
\toprule
Design & Effective sample & EP$\times$green & $p$ \\
\midrule
Full panel (firm level) & 21.5M & $-0.0023$ & 0.57 \\
Collapsed panel & 3.7M cells & $-0.0046$ & 0.51 \\
C-prod-HS4 (within green HS4 families) & 3.8M & $-0.0009$ & 0.86 \\
CEM-matched destinations & 13.7M & $-0.0023$ & 0.54 \\
Deep vs.\ shallow (PTA partners only) & 5.3M & $-0.0022$ & 0.53 \\
Common support (C-overlap) & 21.5M & $-0.0022$ & 0.57 \\
Full panel, with controls & 19.3M & $-0.0002$ & 0.94 \\
Full panel, excluding ASEAN & 19.2M & $-0.0025$ & 0.48 \\
Full panel, including HK--Macao & 23.6M & $-0.0009$ & 0.80 \\
\bottomrule
\end{tabular}
\begin{tablenotes}\footnotesize
\item Each row estimates equation~\eqref{eq:main} on the indicated subsample;
destination-clustered s.e. CEM match on log GDP p.c., growth, MFN tariff (year 2000).
Deep/shallow split at the median of maximum EP depth (16 deep vs.\ 9 shallow countries
in the estimation sample, which excludes Hong Kong and Macao).
The common-support restriction is reported for completeness rather than as an
independent test: 98.5\% of HS6 products are exported both to treated and to
never-treated destinations, so it removes 314 of 21.5 million observations and reproduces
the full panel almost exactly. Its value is in showing that extrapolation outside common
support is not an issue here, not in providing a distinct estimate.
TREND-index analogues (where estimated) give the same picture: full panel $-0.0001$
($p=0.91$), common support $-0.0001$ ($p=0.91$), deep vs.\ shallow $-0.0004$ ($p=0.72$).
Controls: MFN tariff, market concentration (HHI), antidumping indicator.
\end{tablenotes}
\end{threeparttable}
